% ============================================================================
% 02-introducao.tex — Introdução (módulo 3/8)
% ============================================================================
\chapter{Introdução}

\section{Contexto e problema}

O avanço de modelos de linguagem de grande escala (LLMs) transformou o
raciocínio em linguagem natural e matemática: técnicas de raciocínio como
\textit{Chain-of-Thought} \cite{wei_cot}, agentes reflexivos
\cite{shinn_reflexion} e integração com ferramentas \cite{yao_react}
demonstraram ganhos em tarefas de raciocínio. No domínio matemático, a
solução automatizada de problemas da International Mathematical Olympiad
(IMO) tornou-se um marco: o sistema AlphaGeometry \cite{alpha_geometry}
provou 25 de 30 problemas de geometria do nível IMO-2000-2021 e, em 2024, a
combinação AlphaProof/AlphaGeometry alcançou nível de medalha de prata.
Paralelamente, sistemas multiagentes como AutoGen \cite{wu_autogen} e
MetaGPT \cite{hong_metagpt} propuseram arquiteturas de colaboração entre
agentes para tarefas complexas.

Contudo, a avaliação de LLMs em tarefas matemáticas \cite{hendrycks_math,
cobbe_gsm8k, hendrycks_mmlu} raramente é conduzida com \textit{modelos
gratuitos reais}, com \textit{código auditável}, \textit{problemas de olimpíada
não publicados em treino de forma garantida} e \textit{orquestração
multiagente explícita}. Além disso, o campo sofre de alegações infladas:
a literatura alerta que \textit{emergência} aparente pode ser artefato de
métrica \cite{schaeffer_mirage}, e a política anti-\textit{overclaim} do
OpenCode Ecosystem Core (documentada em CORRIGENDUM.md) proíbe classificar
qualquer resultado como \textit{superhuman}, \textit{verificado} ou
\textit{Qualis A1} sem validação externa.

O problema de pesquisa deste estudo é: \textit{como orquestrar LLMs gratuitas
reais, via arquitetura multiagente auditável, para raciocínio matemático de
alto rigor (problemas da IMO), e como produzir evidência empírica honesta —
incluindo limitações — para classificar comparativamente tais modelos?}

\section{Objetivos}

\begin{itemize}
  \item \textbf{Objetivo geral}: medir, com metodologia auditável, o
  raciocínio matemático de LLMs gratuitas reais em problemas canônicos da
  IMO e estruturar o resultado como artigo científico em módulos LaTeX no
  padrão ABNT atualizado, com citações de DOIs ativos.
  \item \textbf{Objetivos específicos}:
  \begin{enumerate}
    \item Curar a paisagem acadêmica (20 agentes + 5 fontes acadêmicas,
    incluindo \texttt{deepmind/acme}) e validar contrato de pipeline real
    (R497--R498).
    \item Implementar solucionador determinístico anti-\textit{leak} para o
    \textit{harness} IMO (R499).
    \item Benchmark de modelos free (\texttt{mimo}, \texttt{big-pickle},
    \texttt{nemotron}, \texttt{muse-spark}) sob condições C0/C1/C2,
    estimando \textit{score} composto.
    \item Replicar o benchmark com roteamento automático (R500) e reportar o
    não-determinismo como limitação.
    \item Produzir fluxogramas de cada processo e raio-x completo da
    arquitetura que conduziu ao resultado.
  \end{enumerate}
\end{itemize}

\section{Justificativa}

A relevância reside em três frentes: (i) \textbf{científica} — a replicação
expõe a fragilidade de \textit{benchmarks} de LLM de \textit{um trial} por
condição, comum na literatura; (ii) \textbf{tecnológica} — o roteamento
automático R500 torna o \textit{benchmark} uma decisão operacional, não um
relatório estático; (iii) \textbf{editorial} — o entregável é um artigo ABNT
em módulos LaTeX compiláveis, com 15 DOIs validados ativos, atendendo ao
rigor de editoras e bancas de periódicos de alto nível (Qualis A1 \textit{ex
ante}), sem prometer aceitação.

\section{Estrutura do relatório}

A seção 2 apresenta a fundamentação teórica; a seção 3 descreve a
metodologia, incluindo o fluxograma do pipeline e o desenho experimental; a
seção 4 reporta os resultados (ranks, replicação e achados); a seção 5
discute o raio-x da arquitetura, os agentes de melhor desempenho e as
limitações; a seção 6 conclui.